% !TEX root = ../main.tex

% 中英标题：\chapter{中文标题}[英文标题]
\chapter{基于结构与语义联合建模的检索增强生成方法}[Joint Semantic Modeling-based RAG]

\section{引言}[Introduction]

近年来，检索增强生成通过引入外部知识，
有效缓解了大语言模型在知识更新与事实准确性方面的局限。
然而，在复杂多跳问答任务中，
如何在保证检索精度的同时降低上下文噪声，
并有效组织检索到的知识信息，
仍然是当前研究面临的关键挑战

针对这一问题，基于知识图谱的 GraphRAG\cite{graphrag} 方法通过显式建模实体及其关系，构建结构化知识以支持多跳推理。
典型方法如 GraphRAG 通过社区划分与分层摘要实现由粗到细的检索过程，
后续工作（如 GFM-RAG\cite{gfm-rag}、GRetriever\cite{g-retriever} 与 LightRAG\cite{lightrag}）进一步引入图神经网络、
结构优化目标及多级索引机制，
在一定程度上提升了推理能力与检索效率。此外，HippoRAG 及其改进版本 HippoRAG2 从认知机制出发，
结合个性化 PageRank 实现多跳检索，增强了上下文建模能力。

然而，现有研究表明，GraphRAG 方法在实际应用中往往难以稳定优于传统 RAG。
这一问题的根本原因在于自动构建知识图谱的质量难以保证：一方面，关系抽取过程存在误差，导致局部语义关系不准确；
另一方面，缺乏全局一致性约束，使得图结构呈现碎片化特征。这些问题共同引入结构冲突与语义歧义，从而削弱检索结果的有效性。

为降低复杂关系建模带来的不稳定性，
LinearRAG\cite{linearrag} 提出以“线性图”为核心的轻量化建模方式，
通过仅建模实体与文本之间的语义关联，
实现无需复杂关系图构建的高效检索框架。
该方法基于分层语义结构与两阶段检索机制，
在多跳问答任务中取得了较好的性能与效率表现。

尽管 LinearRAG 在效率与鲁棒性方面具有优势，
但其仍存在若干局限。
首先，仅依赖实体级语义信息，
难以覆盖更丰富的语义表达单元，例如短语级概念和描述性语义；
其次，段落表示通常基于整体编码，
忽略句子层级的重要性差异，
限制了段落表示的精细程度；
此外，其上下文构建仍以段落为主，
缺乏对实体、短语、知识三元组等语义与结构信息的统一组织，
从而限制了模型对复杂语义关系的理解能力。

为此，本章在 LinearRAG 框架基础上，
提出一种基于结构与语义联合建模的检索增强生成方法
——\textbf{J}oint \textbf{S}emantic LinearRAG（\textbf{JS-LinearRAG}）。
该方法从语义信息扩展、段落表示优化和结构增强上下文构建三个方面对原有流程进行改进，
主要包括以下三个方面：

（1）语义信息抽取机制：
在语义抽取阶段，引入实体、短语与知识三元组的联合抽取。
其中，实体和短语用于扩展语义覆盖范围，
知识三元组作为辅助结构信息补充语义单元之间的关系表达，
从而增强整体语义表达能力；

（2）基于句子加权的段落表示：
在表示学习阶段，通过对句子嵌入进行加权聚合，
刻画段落内部的语义重要性分布，
从而提升段落表示的精细度；

（3）结构增强的联合语义上下文构建：
在上下文构建阶段，将实体、短语、知识三元组及段落信息进行统一组织，
形成兼具语义覆盖能力与结构表达能力的高密度上下文，
以提升模型的理解与推理能力。

\begin{figure}[htbp]
  \centering
  \includegraphics[width=0.98\textwidth]{JSLinearRAG对比.png}
  \caption{四种 RAG 范式的对比分析}
  \label{fig:JSLinearRAG对比}
\end{figure}

与第三章提出的基于结构分析的 TTRAG 方法相比，
本章方法同样引入知识三元组信息，
但二者的作用方式不同。
第三章侧重于利用知识三元组构建显式结构化推理路径，
以实现对多跳推理证据的逐步组织；
而本章并不构建完整的推理路径，
而是将知识三元组作为辅助结构信息，
与实体、短语、知识三元组和段落共同构成联合语义上下文。
因此，本章方法可以看作是在第三章结构分析思想基础上的轻量化扩展，
通过结构信息与语义信息的协同组织，
实现对检索内容的精细表达与高效利用。
该方法为在轻量化建模框架下提升多跳推理能力提供了一种新的解决思路。

本章组织结构如下：4.2节介绍语义信息抽取方法。4.3节介绍基于句子加权的段落表示方法。
4.4节介绍结构增强的联合语义上下文构建方法。4.5节介绍实验结果与分析。4.6节为本章小结。

\section{语义信息抽取}[Multi-granularity Semantic Extraction]

在原始 LinearRAG 中，语义建模主要依赖实体（Entity）作为核心语义单元，
通过实体链接构建文本与图结构之间的关联。
然而，仅依赖实体级语义信息往往难以充分表达复杂文本中的完整语义结构。
在实际多跳问答任务中，除显式命名实体外，
大量关键知识还以短语（Phrase）或更高层次的事实关系形式存在，
例如概念性表达、属性描述及实体间关系。
因此，仅使用实体级表示容易导致语义覆盖不足，
从而限制检索系统对复杂问题的语义理解能力。

为此，本文提出一种语义信息抽取方法，
在实体抽取基础上，进一步引入短语抽取与知识三元组辅助抽取机制，
构建更加丰富的联合语义信息表示。
其中，实体与短语主要用于增强语义覆盖范围，
知识三元组则作为辅助结构信息，
用于补充语义单元之间的潜在关系表达，
从而实现结构信息与语义信息的协同建模。

\subsection{任务定义}

给定检索得到的文档集合：
\begin{equation}
\mathcal{D}_k = \{d_1, d_2, \dots, d_k\},
\end{equation}

\noindent 其中，$k$ 表示检索文档数量，每个文档 $d_i$ 包含若干文本句子。

本章目标是从文档集合中抽取联合语义信息集合 $\mathcal{U}$，
包括实体集合 $\mathcal{E}$、短语集合 $\mathcal{P}$ 与知识三元组集合 $\mathcal{T}$：
\begin{equation}
\mathcal{U} = \mathcal{E} \cup \mathcal{P} \cup \mathcal{T},
\end{equation}

\noindent 其中：
$\mathcal{E}$ 表示实体集合；
$\mathcal{P}$ 表示短语集合；
$\mathcal{T}$ 表示知识三元组集合。

通过联合建模不同层次的语义信息，
本章旨在同时提升语义覆盖能力与结构表达能力。

\subsection{实体抽取}

对于每个文档 $d_i \in \mathcal{D}_k$，
首先采用命名实体识别模型抽取实体集合：
\begin{equation}
\mathcal{E}_i = \text{NER}(d_i),
\end{equation}

\noindent 其中，$\text{NER}(\cdot)$ 表示实体识别函数。

为提升语义质量，
本文过滤掉数值类、序数类等弱语义实体，
仅保留具有实际知识意义的实体节点。

全局实体集合定义为：
\begin{equation}
\mathcal{E} = \bigcup_{i=1}^{k} \mathcal{E}_i,
\end{equation}

实体抽取为后续语义匹配提供基础节点信息。

\subsection{短语抽取}

为弥补实体级表示不足，
本文进一步引入短语级语义单元。
具体而言，
采用依存句法分析中的名词短语提取方法：
\begin{equation}
\mathcal{P}_i = \text{Chunk}(d_i),
\end{equation}

\noindent 其中，$\text{Chunk}(\cdot)$ 表示短语抽取函数。

短语通常能够表达更完整的局部语义，
例如属性、概念或复合描述，
可有效补充实体无法完整表达的语义信息。

全局短语集合定义为：
\begin{equation}
\mathcal{P} = \bigcup_{i=1}^{k} \mathcal{P}_i,
\end{equation}

\subsection{知识三元组辅助抽取}

为进一步增强联合语义信息中的结构表达能力，
本文在实体与短语抽取基础上，
利用开放式信息抽取方法从文本中抽取知识三元组：
\begin{equation}
\mathcal{T}_i = g(d_i),
\end{equation}

\noindent 其中，
$g(\cdot)$ 表示基于大语言模型提示的知识三元组抽取函数，
$\mathcal{T}_i$ 表示文档 $d_i$ 中抽取得到的知识三元组集合。

每个知识三元组可表示为：
\begin{equation}
t=(h,r,t),
\end{equation}

\noindent 其中，
$h$ 表示头实体或主题短语，
$r$ 表示语义关系，
$t$ 表示尾实体、属性或事实描述。

需要强调的是，
与第三章不同，
本章中的知识三元组并不用于构建显式推理路径或复杂路径搜索，
而仅作为联合语义建模中的辅助结构信息，
用于补充语义单元之间的关系表达，
提升整体知识表示完整性。

全局知识三元组集合定义为：
\begin{equation}
\mathcal{T} = \bigcup_{i=1}^{k} \mathcal{T}_i,
\end{equation}

\subsection{联合语义信息组织}

在获得实体集合 $\mathcal{E}$、短语集合 $\mathcal{P}$ 与知识三元组集合 $\mathcal{T}$ 后，
本文构建统一的联合语义信息集合：
\begin{equation}
\mathcal{U} = \{u_1,u_2,\dots,u_n\}, \quad u_i \in \mathcal{E} \cup \mathcal{P} \cup \mathcal{T},
\end{equation}

\noindent 其中，
$n$ 表示联合语义单元总数。

随后，
本文对联合语义信息进行去重、过滤与相关性筛选，
包括删除重复语义单元、过滤低信息量或弱语义单元以及保留与查询相关性较高的语义信息。
最终形成兼具语义覆盖与结构增强能力的联合语义表示。

相比原始 LinearRAG，
本文方法具有以下优势：

（1）语义覆盖增强：实体与短语共同提升语义表达完整性；

（2）结构信息补充：知识三元组提供辅助关系表达；

（3）信息密度提升：联合语义信息组织提高知识表示效率；

（4）检索性能增强：更丰富的信息单元有助于提升复杂多跳检索能力。

因此，
本文提出的语义信息抽取方法为后续段落表示学习与结构增强上下文构建提供了更加全面的基础。

\section{基于句子加权的段落表示方法}[Sentence-weighted Paragraph Representation]

在传统的检索增强生成方法中，段落通常通过整体编码或简单平均池化（mean pooling）的方式进行表示。
然而，这种方式忽略了段落内部句子在语义表达上的差异性，使得重要信息可能被稀释，从而影响检索精度与后续推理效果。

为此，本章提出一种基于句子加权的段落表示方法，通过在句子层面对语义信息进行建模，并基于句子重要性进行加权聚合，
从而获得更加精细的段落表示。

给定文档集合 $\mathcal{D}_k$，每个文档 $d_i$ 由若干句子组成：
\begin{equation}
d_i = \{s_{i1}, s_{i2}, \dots, s_{i n_i}\},
\end{equation}

\noindent 其中， $s_{ij}$ 表示文档 $d_i$ 中的第 $j$ 个句子，$n_i$ 表示句子数量。本章的目标是构建段落表示向量 $\mathbf{e}_{d_i}$，
使其能够更准确地反映段落的语义结构。

\subsection{句子级语义表示}

首先，对每个句子进行编码，得到其向量表示：
\begin{equation}
\mathbf{e}_{s_{ij}} = \text{Encoder}(s_{ij}),
\end{equation}

\noindent 其中，$\mathbf{e}_{s_{ij}}$ 表示文档 $d_i$ 中第 $j$ 个句子 $s_{ij}$ 的语义向量表示，$\text{Encoder}(\cdot)$ 表示文本编码模型（如 E5、BGE\cite{chen2024bge} 等）。相比直接对整个段落编码，
句子级编码能够保留更细粒度的语义信息。

\subsection{句子权重建模}

考虑到不同句子在段落中的信息量存在差异，本章采用基于句子长度的权重分配策略。具体地，将句子的词数作为其重要性度量：
\begin{equation}
w_{ij} = \frac{|s_{ij}|}{\sum_{k=1}^{n_i} |s_{ik}|},
\end{equation}

\noindent 其中， $|s_{ij}|$ 表示句子 $s_{ij}$ 的词数，$w_{ij}$ 表示句子 $s_{ij}$ 在段落 $d_i$ 中的归一化权重，$n_i$ 表示文档 $d_i$ 中的句子总数。

该设计基于如下直观假设：较长的句子通常包含更多语义信息，因此在段落表示中应赋予更高权重。

\subsection{段落表示构建}

在获得句子表示与对应权重后，本章通过加权聚合的方式构建段落表示：
\begin{equation}
\mathbf{e}_{d_i} = \sum_{j=1}^{n_i} w_{ij} \cdot \mathbf{e}_{s_{ij}},
\end{equation}

\noindent 其中，$\mathbf{e}_{d_i}$ 表示文档 $d_i$ 的最终段落向量表示，$w_{ij}$ 表示句子权重，$\mathbf{e}_{s_{ij}}$ 表示对应句子的语义向量。

该表示方式能够根据句子的重要性对语义进行重构，从而避免简单平均带来的信息稀释问题。

与传统的段落表示方法相比，本章提出的句子加权方法具有以下优势：

（1）细粒度语义建模：通过句子级编码保留段落内部的语义结构；

（2）重要性建模：基于句子长度分配权重，使关键信息获得更高贡献；

（3）鲁棒性提升：减少无关或冗余句子对整体表示的干扰。

此外，基于句子加权的段落表示方法无需引入额外训练参数，计算开销较低，易于集成到现有检索框架中。
基于句子加权的段落表示方法能够在保证效率的前提下显著提升语义表示质量，为后续的语义检索与上下文构建提供更加精确的基础。

\section{结构增强的联合语义上下文构建方法}[Joint Semantic Context Construction]

在完成语义信息抽取（第4.2节）与基于句子加权的段落表示（第4.3节）后，
如何将不同类型的语义信息与结构信息进行有效融合，以构建高质量上下文，
是影响检索增强生成性能的关键因素。

传统 RAG 方法通常仅使用段落级文本作为上下文，而 LinearRAG 虽引入实体信息，
但其上下文仍以文本段落这一单一语义来源为主，
缺乏对实体、短语及知识三元组等补充语义与结构信息的统一建模。
这种较为单一的上下文构建方式，限制了模型对复杂语义关系与潜在结构关联的理解能力。

为此，本章提出一种结构增强的联合语义上下文构建方法，
将实体、短语、知识三元组与段落四类信息进行统一组织，
形成兼具语义覆盖能力与结构表达能力的高密度上下文表示。
其中，实体和短语用于增强语义表达范围，
知识三元组用于补充语义单元之间的关系结构，
段落文本则提供完整的上下文证据。

值得注意的是，本章方法在设计思想上受到第三章 TTRAG 方法的启发。
TTRAG 通过构建显式的知识三元组推理路径，
并结合文本证据进行上下文证据索引，
实现了结构化知识与非结构化文本的有效融合。
与之不同，本章并不显式构建推理路径，
而是将知识三元组作为辅助结构信息整合进联合语义上下文中，
通过实体、短语、知识三元组与段落的协同组织，
在隐式语义空间中引导模型完成推理。
因此，本章方法可以看作是在第三章结构分析思想基础上的轻量化语义扩展，
在降低结构推理构建成本的同时，保留一定的结构增强能力。

\subsection{语义单元检索}

基于前述构建的联合语义信息集合
$\mathcal{U} = \mathcal{E} \cup \mathcal{P} \cup \mathcal{T}$
与段落集合 $\mathcal{D}_k$，
首先进行语义检索，得到与查询 $q$ 最相关的语义单元集合与段落集合：
\begin{equation}
\mathcal{U}_q = \text{TopK}(\text{Sim}(q, \mathcal{U})),
\end{equation}

\noindent 其中，
$\mathcal{U}_q$ 表示与查询 $q$ 最相关的前 $K$ 个语义单元集合，
$\text{Sim}(\cdot)$ 表示基于向量表示的相似度函数，
$\text{TopK}(\cdot)$ 表示选择相似度最高的前 $K$ 个结果。

进一步地，段落集合检索定义为：
\begin{equation}
\mathcal{D}_q = \text{TopK}(\text{Sim}(q, \mathcal{D}_k)),
\end{equation}

\noindent 其中，
$\mathcal{D}_q$ 表示与查询 $q$ 最相关的前 $K$ 个段落集合，
$\mathcal{D}_k$ 表示初始检索得到的候选段落集合。

由于 $\mathcal{U}$ 同时包含实体、短语和知识三元组，
因此可将检索结果进一步划分为：
\begin{equation}
\mathcal{U}_q =
\mathcal{U}_q^{(E)} \cup
\mathcal{U}_q^{(P)} \cup
\mathcal{U}_q^{(T)},
\end{equation}

\noindent 其中，
$\mathcal{U}_q^{(E)}$ 表示与查询相关的实体集合，
$\mathcal{U}_q^{(P)}$ 表示与查询相关的短语集合，
$\mathcal{U}_q^{(T)}$ 表示与查询相关的知识三元组集合。

\subsection{知识三元组结构增强}

知识三元组能够显式表达语义单元之间的关系，
因此可作为实体和短语信息的结构补充。
在本章中，
知识三元组并不用于构建完整的显式推理路径，
也不执行第三章中的迭代式路径扩展和束搜索过程，
而是作为结构增强信息参与联合语义上下文构建。

具体而言，
知识三元组 $t=(h,r,o)$ 可提供比单一实体或短语更明确的关系表达。
其中，
$h$ 表示头实体或主题短语，
$r$ 表示语义关系，
$o$ 表示尾实体、属性或事实描述。
通过引入 $\mathcal{T}_q$，
模型能够在阅读段落文本之前先获得简洁的结构化关系提示，
从而提高上下文的信息密度。

需要强调的是，
本节中的知识三元组结构增强主要服务于联合语义上下文组织，
其目标并不是替代第三章中的结构化推理路径，
而是在较低成本下补充实体与短语之间的关系信息。
因此，知识三元组在本章中承担的是辅助结构增强作用，
而非完整推理主体。

\subsection{联合语义上下文构建}

在获得相关实体集合、短语集合、知识三元组集合与段落集合后，
本文将其进行统一组织，构建最终的联合语义上下文：
\begin{equation}
\mathcal{C}_q =
\mathcal{U}_q^{(E)}
\oplus
\mathcal{U}_q^{(P)}
\oplus
\mathcal{T}_q
\oplus
\mathcal{D}_q,
\end{equation}

\noindent 其中，
$\mathcal{C}_q$ 表示最终构建的查询上下文，
$\oplus$ 表示上下文融合操作。

具体地，联合语义上下文由四部分组成：
实体信息、短语信息、知识三元组信息与段落信息。
实体信息负责提供核心语义节点；
短语信息负责补充概念性和描述性语义；
知识三元组信息负责提供辅助关系结构；
段落信息负责提供完整文本证据。
该结构能够在保持语义完整性的同时，
进一步提升上下文的结构性与信息密度。

为进一步提升上下文质量，
本文采用层级化组织策略，
将不同信息按照“语义单元—结构关系—文本证据”的顺序进行排列：
\begin{equation}
\mathcal{C}_q =
\text{Concat}
(
\mathcal{U}_q^{(E)},
\mathcal{U}_q^{(P)},
\mathcal{T}_q,
\mathcal{D}_q
),
\end{equation}

\noindent 其中，
$\text{Concat}(\cdot)$ 表示上下文拼接函数。

该“实体—短语—知识三元组—段落”的层级结构，
有助于引导大语言模型先理解核心语义单元，
再获得简洁的结构关系提示，
最后结合完整文本段落进行综合推理。
如\figref{fig:js_prompt_template}所示，
结构增强的联合语义上下文通过多类型信息组织，
为大语言模型提供了高密度且结构清晰的输入，
从而提升推理能力与答案准确性。

\begin{figure}[htbp]
  \centering
  \begin{minipage}{0.95\linewidth}
      \vspace{2mm}
      \begin{tcolorbox}[title=结构增强联合语义上下文提示模板, fonttitle=\small, fontupper=\small]
          % \small
  
          作为一名高级阅读理解助手，你的任务是仔细分析文本段落和相应的问题。
          你的回答应从“思考：”之后开始，在此部分，你需要系统地阐述你的推理过程，说明你是如何得出结论的。
          最后，以“答案：”结尾，给出简洁明了、明确的回答，无需赘述。
  
  \vspace{5pt}
  
  \texttt{In-Context}
  
  \vspace{3pt}
  
  实体: \{Retrieved Entities\}
  
  \vspace{3pt}
  
  短语: \{Retrieved Phrases\}
  
  \vspace{3pt}

  知识三元组: \{Retrieved Knowledge Triples\}
  
  \vspace{3pt}
  
  文本段落: \{Retrieved Paragraphs\}
  
  \vspace{5pt}

  问题: \{Question\}
  
  \vspace{5pt}
  
  思考：
  
      \end{tcolorbox}
      \vspace{2mm}
  \end{minipage}
  \caption{结构增强联合语义上下文提示模板}
  \label{fig:js_prompt_template}
\end{figure}

相比传统方法，本章提出的结构增强联合语义上下文构建方法具有以下优势：
首先，同时利用实体、短语、知识三元组与段落信息，
能够提升语义覆盖范围与结构表达能力；
其次，知识三元组作为简洁的关系表达形式，
可在较少 Token 开销下为上下文提供结构化提示；
再次，层级化组织不同类型的信息，
有助于大语言模型先把握核心语义与关系线索，
再结合完整段落文本进行综合推理。

与第三章 TTRAG 方法相比，
本章方法不依赖显式构建的完整推理路径，
而是通过结构信息与语义信息的联合建模，
引导模型在隐式语义空间中完成推理。
这种方式在保持较低构建成本的同时，
提供了更加灵活的上下文表达形式。
因此，结构增强联合语义上下文构建方法实现了从“单一语义表示”
向“结构与语义联合建模”的转变，
为提升检索增强生成在复杂问答任务中的性能提供了关键支持。

\section{实验结果及分析}[Experimental Results and Analysis]

\subsection{数据集}

本章在三个典型的多跳问答基准数据集上对所提出方法进行评估，
包括 MuSiQue、2WikiMultiHopQA 和 HotpotQA。实验设置遵循 HippoRAG2 的配置，
从每个数据集的验证集中随机采样 1000 条查询进行评估，所有查询均涉及多跳推理任务。

\subsection{实验设置及评价指标}

在实现方面，本文采用 spaCy 作为实体与短语抽取工具，编码模型选用 all-mpnet-base-v2，用于生成语义向量表示。
知识三元组抽取与最终回答由大语言模型 GPT-4o-mini 或 Llama-3-8B-Instruct 生成。在超参数设置方面，检索文档数量设为 top-$k=5$。

本章采用两个指标对方法进行评估：Contain-Acc.（包含匹配准确率）和GPT-Acc.（基于大语言模型的评估准确率）。

\myparagraph{Contain-Acc}
用于衡量生成答案是否包含标准答案，是一种宽松匹配指标。对于给定问题，若预测答案 $A_{pred}$ 中包含任一标准答案 $A_{true}$，
则视为预测正确。其定义如下：
\begin{equation}
\text{Contain-Acc.} = \mathbb{I}(A_{true} \subseteq A_{pred}),
\end{equation}

\noindent 其中，$\text{Contain-Acc.}$ 表示包含准确率指标，$\mathbb{I}(\cdot)$ 为指示函数，当条件成立时取值为 1，否则为 0。最终结果为所有样本上的平均值。

\myparagraph{GPT-Acc}
是一种基于大语言模型的评估指标，用于衡量预测答案与标准答案在语义层面的匹配程度。
具体而言，采用 GPT-4o-mini 或 Llama3-8B-Instruct 作为评估器，对预测答案 $A_{pred}$ 与标准答案 $A_{true}$ 进行比较，并输出判断结果：
\begin{equation}
\text{GPT-Acc.} = \frac{1}{N} \sum_{i=1}^{N} \text{LLM-Eval}(A_{pred}^{(i)}, A_{true}^{(i)}),
\end{equation}

\noindent 其中，$\text{GPT-Acc.}$ 表示基于大语言模型的准确率指标，$\text{LLM-Eval}(\cdot)$ 表示由大语言模型返回的二值判断结果（匹配为 1，不匹配为 0），$N$ 为测试样本数量，$A_{pred}^{(i)}$ 与 $A_{true}^{(i)}$ 分别表示第 $i$ 个测试样本的预测答案与标准答案。

相比于 Contain-Acc.，GPT-Acc. 能够更好地评估语义一致性，尤其适用于表达方式不完全一致但语义等价的情况。

\subsection{基线模型}

本章将所提出的 JS-LinearRAG 方法与多种具有代表性的 RAG 方法进行对比，这些方法覆盖了从传统密集检索到图结构增强检索的不同技术路线。

\begin{itemize}

\item \textbf{Dense Retrieval (Naive RAG)}
基于向量相似度的经典检索增强生成方法。
该方法直接将查询与文档编码为向量，通过相似度检索相关文本，并将其拼接作为上下文输入大语言模型进行回答，不涉及结构化建模。

\item \textbf{LightRAG}\cite{lightrag}
一种轻量级图结构 RAG 。

\item \textbf{HippoRAG}\cite{hipporag}
受人类记忆机制启发的检索框架，将非结构化文本转化为知识图谱，并通过个性化 PageRank 算法进行多跳检索。

\item \textbf{GFM-RAG}\cite{gfm-rag}
一种基于图结构优化的 RAG 方法，通过引入图神经网络与结构约束，对检索过程进行建模与优化，从而提升多跳推理的准确性与泛化能力。

\item \textbf{HippoRAG2}\cite{hipporag2}
在 HippoRAG 原有图结构基础上引入段落节点及“包含关系”边，同时融合语义相似度与结构信息进行联合检索。

\end{itemize}

\subsection{性能实验结果及分析}

主要实验结果如表\ref{table:js_main_result} 所示。我们在三个多跳问答数据集（MuSiQue、2WikiMultiHopQA 和 HotpotQA）上，
对所提出的 JS-LinearRAG 方法与多种主流 RAG 方法进行了系统对比评估。其中“C-Acc”表示 Contain-Acc.，“G-Acc.” 表示 GPT-Acc.，
而“Avg” 列则为两者的平均值。此外，在表格最后一列，我们还计算了三个数据集上 Contain-Acc. 和 GPT-Acc. 的平均值，
以提供整体性能的综合评估。

\begin{table}[htbp]
  \centering
  \caption{JS-LinearRAG 与多种 RAG 方法在三个多跳问答数据集上的性能对比}
  \label{table:js_main_result}
  \vspace{0.5em}\centering\wuhao

  \resizebox{\textwidth}{!}{
    \begin{tabular}{@{}lcccccccccccc@{}}
    \toprule[1.5pt]
    \multirow{2}{*}{Method} 
    & \multicolumn{3}{c}{MuSiQue} 
    & \multicolumn{3}{c}{2WikiMultiHopQA} 
    & \multicolumn{3}{c}{HotpotQA} 
    & \multicolumn{3}{c}{Average} \\
    \cmidrule(lr){2-4} \cmidrule(lr){5-7} \cmidrule(lr){8-10} \cmidrule(lr){11-13}
    & C-Acc. & G-Acc. & Avg 
    & C-Acc. & G-Acc. & Avg 
    & C-Acc. & G-Acc. & Avg 
    & C-Acc. & G-Acc. & Avg \\
    \midrule[1pt]
    \multicolumn{13}{c}{\textbf{Llama-3-8B-Instruct}} \\ \midrule[1pt]
    TTRAG$_{\text{sk+docs}}$ 
    & \underline{31.10} & \textbf{23.00} & \textbf{27.05} 
    & \textbf{59.70} & \underline{48.50} & \underline{54.10} 
    & \underline{61.80} & \underline{65.30} & \underline{63.55} 
    & \underline{50.87} & \underline{45.60} & \underline{48.23} \\

    TTRAG$_{\text{sk+sens}}$ 
    & \textbf{31.30} & \underline{22.10} & \underline{26.70} 
    & \underline{59.40} & \textbf{49.70} & \textbf{54.55} 
    & \textbf{63.60} & \textbf{66.00} & \textbf{64.80} 
    & \textbf{51.43} & \textbf{45.93} & \textbf{48.68} \\

    JS-LinearRAG 
    & 21.80 & 18.60 & 20.20 
    & 54.70 & 35.00 & 44.85 
    & 59.10 & 50.20 & 54.65 
    & 45.20 & 34.60 & 39.90 \\
    \midrule[1pt]

    \multicolumn{13}{c}{\textbf{GPT-4o-mini}} \\  \midrule[1pt]
    DenseRetrieval 
    & 25.90 & 29.10 & 27.50 
    & 43.80 & 38.60 & 41.20 
    & 52.90 & 54.80 & 53.85 
    & 40.87 & 40.83 & 40.85 \\
    
    LightRAG 
    & 27.40 & 28.60 & 28.00 
    & 55.20 & 39.00 & 47.10 
    & 60.30 & 59.50 & 59.90 
    & 47.63 & 42.37 & 45.00 \\
    
    HippoRAG 
    & 29.30 & 24.10 & 26.70 
    & 66.10 & 59.90 & 63.00 
    & 57.00 & 59.30 & 58.15 
    & 50.80 & 47.77 & 49.28 \\
    
    GFM-RAG 
    & 29.90 & 34.60 & 32.25 
    & 66.80 & 59.60 & 63.20 
    & 62.70 & 65.60 & 64.15 
    & 53.13 & 53.27 & 53.20 \\
    
    HippoRAG2 
    & 31.00 & 35.00 & 33.00 
    & 62.70 & 55.00 & 58.85 
    & 62.90 & 64.30 & 63.60 
    & 52.20 & 51.43 & 51.82 \\
    
    LinearRAG 
    & \underline{31.50} & 34.90 & 33.20 
    & \underline{69.80} & \textbf{64.20} & \textbf{67.00}
    & 67.00 & 66.90 & 66.95 
    & 56.10 & 55.33 & 55.72 \\
    
    TTRAG$_{\text{sk+docs}}$ 
    & 29.90 & \textbf{41.10} & \underline{35.50} 
    & 68.90 & 59.30 & \underline{64.10} 
    & \textbf{71.40} & \textbf{70.90} & \textbf{71.15} 
    & \underline{56.73} & \underline{57.10} & \underline{56.92} \\
    
    TTRAG$_{\text{sk+sens}}$ 
    & 29.20 & \underline{40.70} & 34.95 
    & 66.90 & 56.60 & 61.75 
    & \textbf{71.40} & \textbf{70.90} & \textbf{71.15} 
    & 55.83 & 56.07 & 55.95 \\
    
    JS-LinearRAG 
    & \textbf{35.20} & 40.30 & \textbf{37.75} 
    & \textbf{70.00} & \underline{64.00} & \textbf{67.00} 
    & \underline{70.50} & \underline{69.50} & \underline{70.00} 
    & \textbf{58.57} & \textbf{57.93} & \textbf{58.25} \\
    
    \bottomrule[1.5pt]
    \end{tabular}
  }
\end{table}

\myparagraph{整体性能分析}
从整体结果来看，JS-LinearRAG 在所有方法中取得了最佳的综合性能。在三数据集平均指标上，
JS-LinearRAG 的 Contain-Acc. 和 GPT-Acc. 分别达到 58.57\% 和 57.93\%，
显著优于 Dense Retrieval、LightRAG 以及一系列基于图结构的 RAG 方法。
同时，在所有数据集的 Contain-Acc. 和 GPT-Acc. 指标上，JS-LinearRAG 也表现出稳定优势，
说明该方法不仅能够提升答案的匹配率，同时在语义一致性方面也具有更强能力。

与 HippoRAG、GFM-RAG 及 HippoRAG2 等基于知识图谱的方法相比，JS-LinearRAG 在整体性能上同样取得优势。
例如，相比当前较强的基线 HippoRAG2，JS-LinearRAG 在平均 Contain-Acc. 上提升约 6.37 个百分点。
这一结果表明，在不依赖显式关系建模的情况下，通过联合语义融合，同样能够实现高质量的多跳推理。
这进一步说明，相比复杂的图结构建模，轻量化的语义建模方式在实际应用中具有更好的鲁棒性与效率。

与直接基于实体建模的 LinearRAG 相比，JS-LinearRAG 在整体性能上进一步提升。
其中平均 Contain-Acc. 从 56.10\% 提升至 58.57\%，GPT-Acc. 从 55.33\% 提升至 57.93\%。
该提升验证了联合语义建模的有效性：通过引入短语级语义单元并结合句子加权的段落表示，
能够更全面地刻画语义信息，从而提升检索质量与生成效果。
此外，在 MuSiQue 数据集上，JS-LinearRAG 的提升尤为明显（Avg. 达到 37.75\%），
说明在复杂多跳推理场景下，联合语义信息对推理过程具有更强的支撑作用。

\myparagraph{与 TTRAG 的对比分析}
本文进一步将所提出的 JS-LinearRAG 方法与第三章提出的 TTRAG 方法进行对比分析。
TTRAG 通过构建显式的知识三元组推理路径来引导文本检索与上下文构建，从而为多跳推理提供结构化支持。
需要指出的是，两种方法在检索文档数量设置上存在差异：第三章中 TTRAG 采用 top-$k=10$，
而本章实验统一采用 top-$k=5$。为更公平的比较两种方法，本文也进行了 top-$k=10$ 检索设置下的性能对比实验，
结果如表~\ref{table:top10_comparison} 所示。

\begin{table}[htbp]
  \centering
  \caption{JS-LinearRAG 与 TTRAG 在 top-$k=10$ 检索设置下的性能对比}
  \label{table:top10_comparison}
  \vspace{0.5em}\centering\wuhao

  \resizebox{\textwidth}{!}{
    \begin{tabular}{@{}lcccccccccccc@{}}
    \toprule[1.5pt]
    \multirow{2}{*}{Method} 
    & \multicolumn{3}{c}{MuSiQue} 
    & \multicolumn{3}{c}{2WikiMultiHopQA} 
    & \multicolumn{3}{c}{HotpotQA} 
    & \multicolumn{3}{c}{Average} \\
    \cmidrule(lr){2-4} \cmidrule(lr){5-7} \cmidrule(lr){8-10} \cmidrule(lr){11-13}
    & C-Acc. & G-Acc. & Avg 
    & C-Acc. & G-Acc. & Avg 
    & C-Acc. & G-Acc. & Avg 
    & C-Acc. & G-Acc. & Avg \\
    \midrule[1pt]

    TTRAG$_{\text{sk+docs}}$
    & \underline{31.50} & \textbf{44.00} & \textbf{37.75} 
    & \underline{69.30} & \underline{61.00} & \underline{65.15} 
    & \underline{73.30} & \underline{72.30} & \underline{72.80} 
    & \textbf{58.03} & \textbf{59.10} & \textbf{58.57} \\
    TTRAG$_{\text{sk+sens}}$ 
    & 30.20 & \underline{42.90} & \underline{36.55} 
    & 67.50 & 59.80 & 63.65 
    & \textbf{73.60} & \textbf{72.40} & \textbf{73.00} 
    & 57.10 & \underline{58.37} & \underline{57.73} \\
    JS-LinearRAG
    & \textbf{34.20} & 38.20 & 36.20 
    & \textbf{69.90} & \textbf{63.40} & \textbf{66.65} 
    & 68.00 & 69.00 & 68.50 
    & \underline{57.37} & 56.87 & 57.12 \\

    \bottomrule[1.5pt]
    \end{tabular}
  }
\end{table}

在初始实验中采用不同的检索规模，主要源于两种方法在上下文组织机制上的本质差异。
TTRAG 通过显式构建知识三元组推理路径，对检索得到的候选文档进行结构化筛选和压缩，可显著降低最终输入大语言模型的文本长度，
因此能够在保持提示词规模可控的前提下使用更大的检索范围（top-$k=10$），以充分获取潜在多跳推理证据，能够更好的体现方法的优势。
相比之下，JS-LinearRAG 主要依赖结构与语义联合建模提升信息密度，其上下文虽经过语义优化，但仍保留较多原始段落内容，
因此增加检索规模会导致上下文长度过长而引入更多噪声（这也是TTRAG的优势所在），从而导致性能下降，如表~\ref{table:top10_comparison} 所示。
因此采用领域通常设置的较小的检索规模（top-$k=5$=5）。

从实验结果来看，在 top-$k=5$设置下，JS-LinearRAG 在整体平均指标上取得了最好的性能。
TTRAG$_{\text{sk+docs}}$ 与 TTRAG$_{\text{sk+sens}}$ 在部分数据集上仍表现出较强性能，
尤其是在 HotpotQA 上取得了较高的 Contain-Acc. 与 GPT-Acc.，这表明显式推理路径在复杂语义推理任务中具有一定优势。
然而，JS-LinearRAG 在整体平均指标上仍然优于 TTRAG 方法。这一结果表明，尽管显式推理路径能够提供清晰的推理结构，
但其构建过程依赖于知识三元组质量且计算成本较高；相比之下，JS-LinearRAG 方法通过结构与语义联合建模，在无需显式构建推理路径的情况下，
实现了更高效且稳定的性能提升。

而在 top-$k=10$设置下，在较大的检索规模下，TTRAG$_{\text{sk+docs}}$ 的整体平均性能达到 58.57，
TTRAG$_{\text{sk+sens}}$ 达到 57.73，而 JS-LinearRAG 为 57.12。
实验结果表明，当候选文档数量增加时，TTRAG 通过显式知识三元组推理路径对文本证据进行组织和筛选，能够更充分地利用额外文档中的潜在推理信息，
从而在整体性能上取得更优表现。

与此同时，JS-LinearRAG 在 top-$k=5$ 设置下已经取得了 58.25 的整体平均性能，接近甚至超过部分 top-$k=10$ 设置下的方法。
这说明隐式语义推理范式虽然不显式构建推理路径，但通过结构化增强的联合语义建模，
能够在较少检索文档条件下提供高密度上下文信息，具有更高的检索效率与更低的上下文构建成本。

综合来看，两种方法性能差异并非简单由模型优劣决定，而与各自方法设计目标密切相关。
显式结构化推理与隐式语义推理具有不同优势：
TTRAG 更适合需要充分利用较大候选文档集合、并强调推理路径可解释性的复杂多跳场景；
JS-LinearRAG 则更适合在较低检索规模和较低构建成本下实现稳定性能提升的场景。
二者并非简单替代关系，而是分别从结构化证据组织和联合语义表达两个角度提升检索增强生成性能。

\myparagraph{Llama-3-8B-Instruct 上的性能对比}
实验结果表明，在较小规模模型条件下，TTRAG$_{\text{sk+docs}}$ 与 TTRAG$_{\text{sk+sens}}$ 性能表现优于 JS-LinearRAG。
具体而言，TTRAG$_{\text{sk+sens}}$ 在整体平均指标上达到 48.68，而 JS-LinearRAG 仅为 39.90，存在较明显性能差距。
这一结果说明，当回答模型推理能力受限时，显式结构化推理路径的重要性进一步提升。
TTRAG 所构建的知识三元组推理路径能够为较弱模型提供更加明确的推理步骤和结构化知识组织，从而有效降低复杂多跳推理任务中的推理负担。
相比之下，JS-LinearRAG 虽然通过结构与语义联合建模提升了上下文知识密度，但其推理过程更多依赖模型自身对联合语义信息的隐式整合能力，
因此在模型参数规模较小时，其优势难以充分发挥。

从方法适用性角度看，该结果进一步验证了两种范式的差异：
对于参数规模较大、语义理解能力较强的模型，JS-LinearRAG 能够凭借高效的语义组织方式在较低成本下取得优异性能；
而对于参数规模较小或推理能力相对有限的模型，TTRAG 的显式结构化推理机制则更具优势，
能够通过外部结构化路径补偿模型自身推理能力不足。因此，显式结构化推理（TTRAG）在低容量模型场景下表现出更强的稳定性与适应性，
而隐式联合语义推理（JS-LinearRAG）则更适合高性能模型下的高效率部署。
这也进一步说明，未来将两类方法进行融合，结合结构化推理路径与联合语义表达能力，
可能成为提升不同模型规模下 RAG 系统综合性能的重要方向。

\myparagraph{指标分析}
从指标角度来看，Contain-Acc. 主要反映答案是否被成功召回，而 GPT-Acc. 更关注语义一致性。
JS-LinearRAG 在两个指标上均取得较高分数，说明该方法不仅提升了检索召回能力，同时也增强了生成答案的语义准确性。
特别地，JS-LinearRAG 在 GPT-Acc. 上的稳定提升表明，联合语义信息能够帮助大语言模型更好地理解上下文，
从而生成更加准确且一致的答案。

综上，JS-LinearRAG 在多个数据集和评价指标上均表现出优于现有方法的性能。实验结果验证了：
联合语义建模能够有效提升语义表达能力；
句子加权的段落表示有助于提高检索精度；
联合语义上下文在无需显式推理路径的情况下，仍能够实现高质量多跳推理。

这些结果充分说明，本文提出的结构与语义联合建模方法在性能与效率之间取得了良好平衡，具有较强的实际应用价值。

\subsection{消融实验结果及分析}

为验证本文各模块设计的有效性，本文在 LinearRAG 基础上逐步引入不同组件，构建多组对比实验，
结果如表~\ref{table:js_ablation} 所示。消融实验主要围绕三方面展开：联合语义信息抽取、句子加权表示以及联合语义上下文构建。
为便于描述不同模块组合，本文采用如下命名方式：
ent 表示实体（entity），phr 表示短语（phrase），sw 表示句子加权表示（sentence-weighted embedding），
js 表示联合语义上下文（joint semantic context），kt 表示知识三元组（knowledge triple）。
其中，kt 表示由开放式信息抽取获得的知识三元组结构信息，
作为联合语义建模中的辅助结构增强模块。

\begin{table}[htbp]
  \centering
  \caption{JS-LinearRAG 各模块消融实验结果}
  \label{table:js_ablation}
  \vspace{0.5em}\centering\wuhao

  \resizebox{\textwidth}{!}{
    \begin{tabular}{@{}lcccccccccccc@{}}
    \toprule[1.5pt]
    \multirow{2}{*}{Method} 
    & \multicolumn{3}{c}{MuSiQue} 
    & \multicolumn{3}{c}{2WikiMultiHopQA} 
    & \multicolumn{3}{c}{HotpotQA} 
    & \multicolumn{3}{c}{Average} \\
    \cmidrule(lr){2-4} \cmidrule(lr){5-7} \cmidrule(lr){8-10} \cmidrule(lr){11-13}
    & C-Acc. & G-Acc. & Avg 
    & C-Acc. & G-Acc. & Avg 
    & C-Acc. & G-Acc. & Avg 
    & C-Acc. & G-Acc. & Avg \\
    \midrule[1pt]
    
    LinearRAG$_{\text{ent}}$ 
    & 31.50 & 34.90 & 33.20 
    & 69.80 & 64.20 & 67.00 
    & 67.00 & 66.90 & 66.95 
    & 56.10 & 55.33 & 55.72 \\
    
    LinearRAG$_{\text{phr}}$
    & 31.80 & 35.40 & 33.60 
    & 68.80 & 61.70 & 65.25 
    & 66.60 & 67.40 & 67.00 
    & 55.73 & 54.83 & 55.28 \\
    
    LinearRAG$_{\text{ent+phr}}$
    & 32.90 & 36.90 & 34.90 
    & 68.20 & \textbf{66.20} & 67.20 
    & 66.60 & 68.50 & 67.55 
    & 55.90 & 57.20 & 56.55 \\
    
    LinearRAG$_{\text{ent+sw}}$
    & 33.90 & 38.90 & 36.40 
    & 69.90 & 64.60 & \underline{67.25} 
    & 66.40 & 67.30 & 66.85 
    & 56.73 & 56.93 & 56.83 \\
    
    LinearRAG$_{\text{ent+js}}$
    & 33.80 & 38.90 & 36.35 
    & \textbf{71.70} & \underline{65.40} & \textbf{68.55} 
    & 66.10 & 68.40 & 67.25 
    & 57.20 & 57.57 & 57.38 \\
    
    LinearRAG$_{\text{ent+phr+js}}$
    & \textbf{35.30} & \underline{39.30} & \underline{37.30} 
    & \underline{70.10} & 64.10 & 67.10 
    & 68.00 & \textbf{70.60} & \underline{69.30} 
    & \underline{57.80} & \textbf{58.00} & \underline{57.90} \\

    JS-LinearRAG
    & 35.20 & \underline{40.30} & \underline{37.75} 
    & 70.00 & 64.00 & 67.00
    & \textbf{70.50} & 69.50 & \textbf{70.00} 
    & \textbf{58.57} & \underline{57.93} & \textbf{58.25} \\

    JS-LinearRAG$_{\text{kt}}$
    & \textbf{36.00} & \textbf{40.50} & \textbf{38.25} 
    & 67.50 & 60.80 & 64.15
    & \underline{68.30} & \underline{70.10} & 69.20 
    & 57.25 & 57.13 & 57.20 \\
    
    \bottomrule[1.5pt]
    \end{tabular}
  }
\end{table}


\myparagraph{联合语义信息抽取的有效性}
首先，对比 LinearRAG$_{\text{ent}}$、LinearRAG$_{\text{phr}}$ 与 LinearRAG$_{\text{ent+phr}}$ 三种设置可以观察到，
相比仅使用单一实体或短语信息，
联合引入实体与短语后整体性能呈现稳定提升趋势。
其中，LinearRAG$_{\text{ent+phr}}$ 的整体平均性能由
LinearRAG$_{\text{ent}}$ 的 55.72\%
提升至 56.55\%，
平均提升约 0.83 个百分点。
在 MuSiQue 数据集上，
Avg 指标由 33.20\% 提升至 34.90\%；
在 HotpotQA 上，
Avg 指标由 66.95\% 提升至 67.55\%；
同时在 2WikiMultiHopQA 上，
GPT-Acc. 达到 66.20\%，为该阶段最佳结果。
这一结果表明，
相比仅依赖实体级语义信息，
引入短语能够有效补充更加完整的概念性和描述性语义表达，
增强语义覆盖范围，
从而提升检索质量与答案生成能力。
因此，第4.2节提出的联合语义信息抽取机制能够为后续检索增强生成提供更加丰富的语义基础。

\myparagraph{句子加权表示的影响}
在语义信息抽取基础上，
进一步引入句子加权段落表示后，
模型性能继续提升。
LinearRAG$_{\text{ent+sw}}$ 的整体平均性能达到 56.83\%，
相比原始 LinearRAG$_{\text{ent}}$ 提升约 1.11 个百分点。
其中，
MuSiQue 数据集上的 Avg 指标由 33.20\% 提升至 36.40\%，
提升最为显著。
这一结果说明，
相较于直接对段落整体编码，
基于句子重要性加权的表示方式能够更精细地建模段落内部语义结构，
突出关键句子信息，
减少冗余内容干扰，
从而提升段落表示质量。
这验证了第4.3节提出的句子加权表示方法在语义表示精度方面的有效性。

\myparagraph{联合语义上下文构建的作用}
进一步引入联合语义上下文构建后，
模型整体性能获得持续提升。
LinearRAG$_{\text{ent+js}}$ 在 2WikiMultiHopQA 数据集上取得 68.55\% 的最高 Avg 指标，
说明仅通过联合语义上下文组织，
即使不引入短语信息，
也能够显著增强信息组织能力。
当进一步结合完整联合语义抽取后，
LinearRAG$_{\text{ent+phr+js}}$ 的整体平均性能达到 57.90\%，
较基础 LinearRAG$_{\text{ent}}$ 提升约 2.18 个百分点。
其中：
MuSiQue Avg 达到 37.30\%，
HotpotQA Avg 达到 69.30\%，
整体 Contain-Acc. 与 GPT-Acc. 分别达到 57.80\% 和 58.00\%。
这一结果表明，
联合语义抽取与联合上下文构建之间存在明显协同作用，
通过将实体、短语、知识三元组与段落进行层级化组织，
能够有效提升上下文知识密度与模型推理能力。
在此基础上，
进一步引入句子加权表示，
即完整模型 JS-LinearRAG
（LinearRAG$_{\text{ent+phr+sw+js}}$），
整体平均性能达到当前最佳，
Avg 达到 58.25\%。
尤其在 HotpotQA 上，
Contain-Acc. 达到 70.50\%，
整体 Avg 达到 70.00\%，
表现最优。
这说明，
句子级表示优化与联合语义上下文构建能够形成更强的互补作用，
共同提升复杂多跳问答任务中的检索与推理表现。

\myparagraph{知识三元组结构增强的影响}
为进一步分析结构信息在第四章联合建模框架中的作用，
本文在完整 JS-LinearRAG 基础上引入知识三元组结构增强模块，
即 JS-LinearRAG$_{\text{kt}}$（LinearRAG$_{\text{ent+phr+sw+js+kt}}$）。
实验结果表明，
知识三元组结构增强在 MuSiQue 数据集上取得局部最优结果，
Avg 达到 38.25\%，
说明知识三元组提供的结构关系信息在部分复杂推理场景中能够进一步增强知识表达能力。
然而，
在 2WikiMultiHopQA 数据集上，
性能出现较明显下降，
整体平均性能由 58.25\% 降至 57.20\%。
这一现象说明，
虽然知识三元组能够提供额外结构关系信息，
但在第四章的轻量化联合语义框架中，
实体与短语已能够较充分覆盖核心语义，
额外引入知识三元组可能带来一定结构冗余与噪声，
从而影响整体稳定性。
因此，
相比第三章中以知识三元组为核心的显式结构化推理路径，
第四章更适合将知识三元组作为辅助增强信息，
而非主要推理主体。
这一结果进一步验证了本文提出的 JS-LinearRAG 的设计合理性：
在轻量级检索增强生成框架下，
优化语义覆盖能力与信息组织方式，
通常比引入更复杂的结构化信息具有更好的综合效率与稳定性。

\subsection{样例分析}

为进一步分析所提出方法在实际多跳推理任务中的表现，本文选取一个典型案例进行对比分析。待回答的问题为：
\textit{“Which film has the director who was born later, Playing It Wild or I'll Be Going Now?”}

如图~\ref{fig:js_case} 所示，左侧为 LinearRAG 的检索结果，右侧为本文提出的 JS-LinearRAG 的检索结果。
可以观察到，LinearRAG 在前 $k=5$ 个检索文档中未能返回与问题直接相关的关键信息，导致后续模型缺乏有效证据支撑，
难以完成正确推理。这一现象表明，当仅依赖实体级语义并采用统一的整体段落嵌入表示进行检索时，
模型难以刻画段落内部的语义差异，容易受到表面语义匹配的干扰，从而遗漏关键的推理路径信息。

相比之下，JS-LinearRAG 能够准确检索到与问题相关的核心实体与语义信息。
具体而言，模型成功定位到两个目标电影 “Playing It Wild” 和 “I'll Be Going Now”，
并进一步检索到其对应导演分别为 “William Duncan” 和 “Dino Risi”。
在此基础上，模型继续获取两位导演的出生日期信息，其中 “William Duncan” 出生于 1879 年 12 月 16 日，
而 “Dino Risi” 出生于 1916 年 12 月 23 日。通过对两者出生时间的比较，
可以正确推断出 “I'll Be Going Now” 的导演出生更晚，从而得到正确答案。

进一步分析 JS-LinearRAG 的优势来源，可以发现其不仅受益于联合语义建模与联合语义上下文构建，
还得益于基于句子加权的段落表示方法。相比于传统的整体段落编码方式，
该方法通过对段落中各句子的嵌入进行加权聚合（以句子长度为权重），
能够更准确地刻画段落内部的语义分布，使包含关键事实信息的句子在表示中占据更大权重。
在本案例中，与导演及其出生时间相关的句子通常信息密度较高且长度较长，因此在加权表示中被显著强化，
从而提升了相关段落在检索阶段的排序位置。

这一机制有效缓解了传统方法中“关键信息被淹没于冗余文本”的问题，使模型能够更加稳定地捕捉跨文档的关键语义线索，
并支持多跳推理过程中的信息衔接。相比之下，LinearRAG 采用统一的段落表示方式，
难以区分段落内部不同句子的语义重要性，从而在复杂查询下更容易出现检索偏差。

\begin{figure}[htbp]
  \centering
  \includegraphics[width=1.0 \linewidth]{js_case.pdf}
  \caption{LinearRAG 与 JS-LinearRAG 的文档检索结果样例}
  \label{fig:js_case}
\end{figure}

\section{本章小结}[Brief summary]

本章针对传统 RAG 方法在复杂多跳问答任务中存在的语义表达不足与上下文利用效率有限等问题，
提出了一种基于结构与语义联合建模的检索增强生成方法（JS-LinearRAG）。
该方法从联合语义表示与融合的角度出发，对 LinearRAG 框架进行了系统性改进。
具体而言，首先在语义抽取阶段引入实体、短语与知识三元组的联合抽取机制，扩展了语义信息的覆盖范围，并利用知识三元组提供辅助结构信息；
其次，在表示学习阶段提出基于句子加权的段落表示方法，通过建模句子级语义分布，提高段落表示的精细程度；
最后，在上下文构建阶段提出结构增强的联合语义上下文构建方法，
将实体、短语、知识三元组与段落信息进行统一组织，
从而提升上下文的信息密度与表达能力。
实验结果表明，所提出方法在多个多跳问答数据集上均取得了优于现有方法的性能表现。
消融实验进一步验证了各模块设计的有效性，
说明结构与语义联合建模在提升检索质量与推理能力方面具有显著作用。
与第三章提出的基于结构分析的 TTRAG 方法相比，
本章方法不依赖显式推理路径，
而是通过结构增强下的联合语义信息隐式建模实现推理过程，
在保证性能的同时具有更高的灵活性与计算效率。
